\documentclass[11pt]{amsart}

\usepackage{amsmath}
\usepackage{amssymb}
\usepackage{amsthm}
\usepackage{booktabs}
\usepackage{array}
\usepackage{alltt}
\usepackage{hyperref}

\newtheorem{theorem}{Theorem}[section]
\newtheorem{proposition}[theorem]{Proposition}
\newtheorem{lemma}[theorem]{Lemma}
\newtheorem{corollary}[theorem]{Corollary}
\theoremstyle{definition}
\newtheorem{definition}[theorem]{Definition}
\theoremstyle{remark}
\newtheorem{remark}[theorem]{Remark}

\numberwithin{equation}{section}
\setlength{\emergencystretch}{2em}
\raggedbottom

\newcommand{\NN}{\mathbb{N}}
\newcommand{\ZZ}{\mathbb{Z}}
\newcommand{\QQ}{\mathbb{Q}}
\newcommand{\RR}{\mathbb{R}}
\newcommand{\CC}{\mathbb{C}}
\newcommand{\cO}{\mathcal{O}}
\newcommand{\Nm}{\operatorname{N}}
\newcommand{\eu}{\phi_{\mathrm E}}
\newcommand{\Tot}{\mathcal T}
\newcommand{\lean}[1]{{\def\_{\textunderscore\allowbreak}\texttt{#1}}}
% Visible marker for open items. Every \TODO must be resolved (or the passage deleted) before posting.
\newcommand{\TODO}[1]{\textbf{[TODO: #1]}}

\begin{document}

\hypersetup{pdftitle={Irrationality of odd q-zeta values at real quadratic units},pdfauthor={Shawn Krisman}}
\title[Odd $q$-zeta values at real quadratic units]{Irrationality of odd $q$-zeta values at real quadratic units}

\author{Shawn Krisman}

\subjclass[2020]{Primary 11J72; Secondary 11B39, 11R11, 33D15, 68V20}
\keywords{$q$-zeta values, Lambert series, real quadratic units, Lucas sequences, Fibonacci numbers, irrationality,
basic hypergeometric series, Lean}

\begin{abstract}
For $|q|<1$ and odd $l\ge3$ let $\zeta_q(l)=\sum_{m\ge1}m^{l-1}q^m/(1-q^m)$. Let $q$ be a unit of norm $+1$ in a
real quadratic field $\QQ(\sqrt d)$ with $0<|q|<1$. We prove that the dimension of the $\QQ$-vector space spanned by
$1,\sqrt d\,\zeta_q(3),\sqrt d\,\zeta_q(5),\dots,\sqrt d\,\zeta_q(A-1)$ is at least $(\pi/\sqrt{12+\pi^2}+o(1))\sqrt A$,
and that one of $\sqrt d\,\zeta_q(5),\dots,\sqrt d\,\zeta_q(11)$ lies outside $\QQ+\QQ\sqrt d\,\zeta_q(3)$; these
seem to be the first results on odd $q$-zeta values of weight at least five at an algebraic irrational $q$. We also give
an elementary proof that $\sqrt d\,\zeta_q(3)$ is irrational, formalised in Lean~4 with Mathlib using only the standard
axioms. In particular, for a Lucas sequence $(U_m)$ of a real quadratic unit $\eta>1$,
$\sum_{m\ge1}m^2(U_{m+1}/U_m-\eta)\notin\QQ$, for instance $\sum_{m\ge1}m^2(F_{m+1}/F_m-\varphi)\notin\QQ$. The proofs
use the very-well-poised basic hypergeometric series of Krattenthaler, Rivoal and Zudilin, whose coefficients turn out
to be invariant or anti-invariant under $q\mapsto1/q$, up to a power of $q$; at a unit with conjugate $1/q$ this makes suitable integer linear
forms in the twisted values $\sqrt d\,\zeta_q(l)$ small at both archimedean places. A stronger form of the weight-three statement has been claimed concurrently and
independently in unverified preprints.
\end{abstract}

\maketitle

% ===========================================================================
\section{Introduction}
% ===========================================================================

For a complex number $q$ with $|q|<1$ and an integer $l\ge1$ let
\[
\zeta_q(l)=\sum_{m\ge1}\frac{m^{l-1}q^m}{1-q^m}=\sum_{N\ge1}\sigma_{l-1}(N)\,q^N,
\]
the $q$-analogue of $\zeta(l)$ studied by Krattenthaler, Rivoal and Zudilin \cite{KRZ06}. For even $l$ these are, up to
constants, Eisenstein series, and Nesterenko's theorem \cite{Nes96} shows that $\zeta_q(2m)$ is transcendental for
every algebraic $q$ with $0<|q|<1$ (see \cite[Introduction]{JM10}). The odd values are much less understood. Let
$q=1/p$ with $p\in\ZZ\setminus\{0,\pm1\}$. For $p\ge2$, $\zeta_q(1)$ is irrational \cite{Bor92}, and
$1,\zeta_q(1),\zeta_q(2)$ are linearly independent over $\QQ$ \cite{PVA07}. For odd $l\ge3$ Krattenthaler, Rivoal and
Zudilin proved that the dimension of the $\QQ$-span of $1,\zeta_q(3),\zeta_q(5),\dots,\zeta_q(A-1)$ grows at least like
a constant times $\sqrt A$, and that one of $\zeta_q(3),\dots,\zeta_q(11)$ is irrational \cite{KRZ06}; Jouhet and
Mosaki reduced this to one of $\zeta_q(3),\dots,\zeta_q(9)$ \cite{JM10}, and Fischler and Zudilin refined these results
\cite{FZ10}. In October 2026 a proof that $\zeta_{1/p}(3)=\sum_{n\ge1}n^2/(p^n-1)$ is irrational for every integer
$p\ge2$ was claimed in an unrefereed, AI-assisted research draft with a Lean formalisation, posted under the pen name
uplift-universal-inc \cite{UPL,QZ3}.

At algebraic irrational $q$, apart from the even weights, the known results are of weight one, for example for Lambert
series $\sum_na_n/(b^n-1)$ with periodic coefficients $a_n$ at suitable algebraic integers $b$ with $|b|>1$ (Tachiya
\cite{Tac04}, Bundschuh and V\"a\"an\"anen \cite{BV05}, Dixit, Kumar and Pathak \cite{DKP24}; see
Section~\ref{sec:related}). Krattenthaler, Rivoal and Zudilin remark \cite[\S3.5]{KRZ06} that their theorems would extend to algebraic
$q$, via Jadot's linear independence criterion over number fields \cite{Jad96}, if a conjectured asymptotic for the
coefficients of their linear forms held. Bundschuh and Zudilin list the arithmetic of $\zeta_q(k)$ at reciprocals of
Pisot and Salem numbers among their open problems \cite[\S6]{BZ07}. Before the concurrent work discussed below, we know
of no result on $\zeta_q(3)$ at an algebraic irrational $q$, not even of ``one of'' type.

A unit $q$ of a real quadratic field $K$ has norm $+1$ exactly when its Galois conjugate is $1/q$. Our results concern
such $q$. Throughout, $K$ is a real quadratic field, $\sigma$ its nontrivial automorphism, $q\in\cO_K^\times$ a unit of
norm $+1$ with $0<|q|<1$, so that $\sigma q=1/q$, and $d\ge1$ an integer with $K=\QQ(\sqrt d)$.

\begin{theorem}\label{thm:many}
For every even $A\ge4$ and every integer $r$ with $1\le r<A/2$,
\[
\dim_\QQ\bigl(\QQ+\QQ\sqrt d\,\zeta_q(3)+\QQ\sqrt d\,\zeta_q(5)+\dots+\QQ\sqrt d\,\zeta_q(A-1)\bigr)\ge1+\tau(A,r),
\]
where
\[
\tau(A,r)=\frac{4\pi^2r(A-2r)-12A}{12A+\pi^2(A+4r^2)}.
\]
In particular, as $A\to\infty$ through even integers, this dimension is at least $(\pi/\sqrt{12+\pi^2}+o(1))\sqrt A
=(0.6717\ldots+o(1))\sqrt A$.
\end{theorem}

\begin{corollary}\label{cor:many}
At least one of $\sqrt d\,\zeta_q(5)$, $\sqrt d\,\zeta_q(7)$, $\sqrt d\,\zeta_q(9)$, $\sqrt d\,\zeta_q(11)$ lies
outside $\QQ+\QQ\sqrt d\,\zeta_q(3)$, and so is irrational. Infinitely many of the numbers $\sqrt d\,\zeta_q(l)$, $l$ odd,
are irrational.
\end{corollary}

The corollary is the case $A=12$, $r=2$, where $1+\tau=2.16\ldots>2$. Theorem~\ref{thm:many} is not formalised. Its case
$A=4$, $r=1$ gives $1+\tau=1.24\ldots>1$, that is, the following theorem, for which we give a separate, shorter proof that
uses only an elementary bound for $\sum_{\ell\le n}\eu(\ell)$ and is formalised in Lean (Section~\ref{sec:lean}).

\begin{theorem}\label{thm:main}
$\sqrt d\,\zeta_q(3)$ is irrational.
\end{theorem}

Equivalently: if $q$ is a real number with $0<|q|<1$ and $s=q+q^{-1}\in\ZZ$, then $(q-q^{-1})\zeta_q(3)\notin\QQ$.
Indeed, such $q$ is a root of $x^2-sx+1$ with $|s|\ge3$, so it is a quadratic unit of norm $1$, and $q-q^{-1}=
\pm\sqrt{s^2-4}$ is a nonzero rational multiple of $\sqrt d$. Neither theorem says anything about the values
$\zeta_q(l)$ themselves; see Remark~\ref{rem:plain}.

For integers $P,Q$ let $U_m=U_m(P,Q)$ be the Lucas sequence $U_0=0$, $U_1=1$, $U_{m+1}=PU_m-QU_{m-1}$. Call $(P,Q)$
\emph{admissible} if $Q=-1$ and $P\ge1$, or $Q=1$ and $P\ge3$. Then $D=P^2-4Q>0$ is not a square, $\eta=(P+\sqrt D)/2$
is a unit of norm $Q$ of the order $\ZZ[\eta]$, $\eta>1$, and every unit $\eta>1$ of a real quadratic order arises in
this way \cite[\S1]{Kri26b}. The Fibonacci numbers are $F_m=U_m(1,-1)$, with $\eta=\varphi=(1+\sqrt5)/2$, and the Pell
numbers are $\mathcal P_m=U_m(2,-1)$, with $\eta=1+\sqrt2$.

\begin{corollary}\label{cor:lucas}
For every admissible $(P,Q)$, with $\eta=(P+\sqrt{P^2-4Q})/2$,
\[
\sum_{m\ge1}m^2\Bigl(\frac{U_{m+1}}{U_m}-\eta\Bigr)\notin\QQ.
\]
\end{corollary}

\begin{corollary}\label{cor:fib}
The numbers
\[
\sum_{m\ge1}m^2\Bigl(\frac{F_{m+1}}{F_m}-\varphi\Bigr)=\sum_{m\ge1}\frac{(-1)^mm^2}{\varphi^mF_m}
\qquad\text{and}\qquad
\sum_{m\ge1}m^2\Bigl(\frac{\mathcal P_{m+1}}{\mathcal P_m}-1-\sqrt2\Bigr)
\]
are irrational.
\end{corollary}

The ratios $U_{m+1}/U_m$ converge to $\eta$ geometrically, so the series converge absolutely. Corollary~\ref{cor:lucas}
is Theorem~\ref{thm:main} at $q=\bar\eta/\eta=Q\eta^{-2}$, where $\bar\eta=Q/\eta$ (Section~\ref{sec:proof}). It
covers only the units of the form $q=\pm\eta^{-2}$, so it is weaker than Theorem~\ref{thm:main}: for example $q=2-\sqrt3$ is not of
this form, since the units of every order of $\QQ(\sqrt3)$ are $\pm(2+\sqrt3)^k$. In the same way,
$\sum_{m\ge1}m^{l-1}(U_{m+1}/U_m-\eta)=\sqrt{P^2-4Q}\,\zeta_q(l)$ for every odd $l\ge3$, so Theorem~\ref{thm:many} and
Corollary~\ref{cor:many} hold for these sums; for instance, one of $\sum_{m\ge1}m^{l-1}(F_{m+1}/F_m-\varphi)$,
$l\in\{5,7,9,11\}$, is irrational.

\subsection*{Concurrent work}
While this work was being completed, the contributor twinheterodox submitted the Hexagon preprint \cite{Twi26}
(submitted 18 November, published 19 November 2026). It is labelled by its author as primarily AI-generated, and its
proof is not formally verified. It claims that $Z(b)=\sum_{n\ge1}n^2/(b^n-1)=\zeta_{1/b}(3)$ does not lie in $\QQ(b)$
for every algebraic integer $b$ with $|b|>1$ all of whose conjugates outside the closed unit disk belong to
$\{b,\bar b\}$; this includes all Pisot and Salem numbers. For $b=1/q$ with $q$ as in Theorem~\ref{thm:main}, this is
the statement $\zeta_q(3)\notin K$, which is stronger than Theorem~\ref{thm:main}: it implies both
$\sqrt d\,\zeta_q(3)\notin\QQ$ and $\zeta_q(3)\notin\QQ$, and the second does not follow from our method. A companion
preprint by the same contributor \cite{Twi26b} claims that $1$ and the three numbers $\sum_{n\ge1}(b^n-1)^{-j}$,
$j=1,2,3$, are linearly independent over $\QQ(b)$ for every algebraic integer $b$ with $|b|>1$ all of whose
conjugates other than $b$ and $\bar b$ lie in the closed unit disk (by its abstract, a slightly different wording of the
same class); since
$\zeta_{1/b}(3)=M_1+3M_2+2M_3$ with $M_j=\sum_{n\ge1}(b^n-1)^{-j}$ (expand $(b^n-1)^{-j}$ in powers of $b^{-n}$ and use
$k^2=2\binom{k-1}2+3(k-1)+1$), this also contains Theorem~\ref{thm:main}, as $b=1/q$ is in that class. The proof in \cite{Twi26} builds on
the moving-denominator Hankel construction of \cite{UPL}. We have not checked their arguments, and our work is independent of
them, and we claim no priority for Theorem~\ref{thm:main}. Relative to \cite{Twi26,Twi26b}, the contributions of the
present paper are:
\begin{itemize}
\item Theorem~\ref{thm:many} and Corollary~\ref{cor:many}, which concern weights above three, where the claims of
\cite{Twi26,Twi26b} say nothing;
\item an independent proof of Theorem~\ref{thm:main}, which is elementary and short: it uses no estimate on the
coefficients of the linear forms and no input from analytic number theory;
\item a complete formal verification of Theorem~\ref{thm:main} and Corollaries~\ref{cor:lucas} and~\ref{cor:fib} in
Lean~4 with Mathlib, using only the standard axioms (Section~\ref{sec:lean});
\item the Lucas-sequence form of these results; for Theorem~\ref{thm:main} it follows equally from \cite{Twi26}.
\end{itemize}

\subsection*{The method}
We use a variant of the very-well-poised basic hypergeometric series of \cite[\S3]{KRZ06} (with weight $T^e$ instead of
$T^{e+1}$, summed bilaterally), in the case $A=4$, $r=1$. For each $n$ it gives a linear form $S_n(q)=P_0(q)+P_3(q)\zeta_q(3)$ with $P_0,P_3\in\QQ(q)$
(Lemma~\ref{lem:expansion}). The new observation is a reciprocity under $q\mapsto1/q$ (Lemma~\ref{lem:recip}):
\[
P_3(1/q)=q^{-n}P_3(q),\qquad P_0(1/q)=-q^{-n}P_0(q).
\]
At a unit with $\sigma q=1/q$, suppose that $\zeta_q(3)\in\QQ\sqrt d$. Then $S_n\in K$, and Galois conjugation gives
$\sigma(S_n)=-q^{-n}S_n$, so the linear form is small at both archimedean places of $K$, and the product formula gives a contradiction: after clearing the cyclotomic
denominators $d_n(q)^4$ (Lemma~\ref{lem:den}), the norm is about
\[
|\Nm d_n(q)|^4\,|S_n(q)|^2\le|q|^{-(\frac32-2+o(1))n^2},
\]
since $\log|S_n|\sim-n^2\log(1/|q|)$ (Lemma~\ref{lem:size}) and $\log|\Nm d_n(q)|\le(\frac38+o(1))n^2\log(1/|q|)$
(Lemma~\ref{lem:cost}). No upper bound on the coefficients $P_0,P_3$ is needed. The same construction gives nothing
for $\zeta_q(3)$ itself (Remark~\ref{rem:plain}), nor for the Fibonacci zeta value $\sum_mF_m^{-3}$
(Remark~\ref{rem:fibzeta}). Our earlier proofs that $\sum_mF_m^{-3}$ and $\sum_mU_m^{-3}$ are irrational
\cite{Kri26a,Kri26b} used a different construction, moving-denominator Hankel determinants.

For Theorem~\ref{thm:many} we take general $(A,r)$, where $S_n=P_0+\sum_lP_l\zeta_q(l)$ involves the odd weights
$3\le l\le A-1$. The same reciprocity shows that the integer forms
$\operatorname{Tr}_{K/\QQ}(\delta P_0)+\sum_l\operatorname{Tr}_{K/\QQ}(\delta P_l/\sqrt d)\,\sqrt d\,\zeta_q(l)$, for a
suitable denominator $\delta\in\cO_K$, equal $(\delta-q^{-e}\sigma\delta)S_n$. Multiplying $\delta$ by a power of the
unit $q$ splits the cost of the denominator evenly between the two real places, a standard balancing step, and
Nesterenko's linear independence criterion then gives Theorem~\ref{thm:many}. Here we need upper bounds for the
coefficients, as in \cite{KRZ06}, and the asymptotic $\sum_{\ell\le n}\eu(\ell)\sim3n^2/\pi^2$. With the same leading
constants (decay, coefficient size and denominator, which our variant of the series shares with that of \cite{KRZ06})
and the same single-form criterion, the bound $1+\tau(A,r)$ is exactly twice the bound that Krattenthaler, Rivoal and
Zudilin obtain at $q=1/p$ (Remark~\ref{rem:balance}); this compares one method at two kinds of $q$, not with the best results known at $q=1/p$, which are better than those of \cite{KRZ06}
\cite{JM10,FZ10}.

Sections~\ref{sec:R}--\ref{sec:size} hold for every even $A\ge4$ and every $1\le r<A/2$; only
Sections~\ref{sec:cost}--\ref{sec:proof} specialise to $A=4$, $r=1$, and Section~\ref{sec:many} returns to general
$(A,r)$. The identities of Section~\ref{sec:R} are those of \cite{KRZ06} in our normalisation; we give proofs to keep
the paper self-contained.

\subsection*{Formal verification}
Theorem~\ref{thm:main}, Corollaries~\ref{cor:lucas} and~\ref{cor:fib}, and every step of their proofs (for $A=4$,
$r=1$) are formalised in Lean~4 with Mathlib (\url{https://github.com/telaviv/qzeta3-quadratic-units}; Section~\ref{sec:lean}).
Theorem~\ref{thm:many} and Corollary~\ref{cor:many} are not formalised.

% ===========================================================================
\section{The rational function and its symmetries}\label{sec:R}
% ===========================================================================

Write $(x;q)_m=\prod_{i=0}^{m-1}(1-xq^i)$. Fix an even integer $A\ge4$ and an integer $r$ with $1\le r<A/2$, and put
$e=(A-2r)n/2\in\ZZ$. For $n\ge1$ let
\begin{equation}\label{eq:R}
R(T)=R_n(T;q)=\frac{(q;q)_n^{A-2r}\,(q^{-rn}T;q)_{rn}\,(q^{n+1}T;q)_{rn}\,T^e}{(T;q)_{n+1}^A},
\end{equation}
a rational function of $T$ over the field $\QQ(q)$, $q$ an indeterminate. This is the rational function $R_n(T;q)$ of
\cite[\S3.3]{KRZ06} without its factor $q^{-(A-2r)n/4}$. Krattenthaler, Rivoal and Zudilin sum $q^kR_n(q^k)$; we sum
$R_n(q^k)$. Its numerator has degree $2rn+e=(A+2r)n/2<A(n+1)$ in $T$, so $R$ is proper.

\begin{lemma}[well-poised symmetry]\label{lem:wp}
$R(q^{-n}/T)=q^{An/2}\,T^A\,R(T)$.
\end{lemma}

\begin{proof}
Substitute $T\mapsto q^{-n}/T$. The factors $(T;q)_{n+1}$, $(q^{-rn}T;q)_{rn}$ and $(q^{n+1}T;q)_{rn}$ become
$\prod_{i=0}^n(1-q^{i-n}/T)$, $\prod_{m=n+1}^{n+rn}(1-q^{-m}/T)$ and $\prod_{m=1}^{rn}(1-q^{m}/T)$. Taking out of each
linear factor its term in $1/T$ gives
\begin{align*}
(T;q)_{n+1}&\mapsto(-1)^{n+1}q^{-n(n+1)/2}\,T^{-(n+1)}\,(T;q)_{n+1},\\
(q^{-rn}T;q)_{rn}&\mapsto(-1)^{rn}q^{-rn^2-rn(rn+1)/2}\,T^{-rn}\,(q^{n+1}T;q)_{rn},\\
(q^{n+1}T;q)_{rn}&\mapsto(-1)^{rn}q^{rn(rn+1)/2}\,T^{-rn}\,(q^{-rn}T;q)_{rn},\\
T^e&\mapsto q^{-ne}\,T^{-e}.
\end{align*}
The two numerator blocks trade places. Relative to $R(T)$, the power of $T$ is $A(n+1)-2rn-2e=A$, the sign is
$(-1)^{2rn+A(n+1)}=1$, and the power of $q$ is $-rn^2-ne+An(n+1)/2=An/2$.
\end{proof}

\begin{lemma}[$q\mapsto1/q$]\label{lem:inv}
$R(T;1/q)=q^{-e}\,R(q^{-n}T;q)$.
\end{lemma}

\begin{proof}
First, $R(1/V;1/q)=q^{rn}V^AR(V;q)$. The computation is of the same kind:
\begin{align*}
(1/q;1/q)_n&=(-1)^nq^{-n(n+1)/2}\,(q;q)_n,\\
(V^{-1};1/q)_{n+1}&=(-1)^{n+1}q^{-n(n+1)/2}\,V^{-(n+1)}\,(V;q)_{n+1},\\
(q^{rn}V^{-1};1/q)_{rn}&=(-1)^{rn}q^{rn(rn+1)/2}\,V^{-rn}\,(q^{-rn}V;q)_{rn},\\
(q^{-n-1}V^{-1};1/q)_{rn}&=(-1)^{rn}q^{-rn^2-rn(rn+1)/2}\,V^{-rn}\,(q^{n+1}V;q)_{rn}.
\end{align*}
The total power of $q$ is $-(A-2r)n(n+1)/2-rn^2+An(n+1)/2=rn$, the power of $V$ is $A$ and the sign is $+1$. By
Lemma~\ref{lem:wp} with $T=q^{-n}V$, $R(1/V;q)=q^{-An/2}V^AR(q^{-n}V;q)$. Hence
$R(V;1/q)=q^{rn}V^{-A}R(1/V;q)=q^{rn-An/2}R(q^{-n}V;q)$, and $rn-An/2=-e$.
\end{proof}

% ===========================================================================
\section{Partial fractions and denominators}\label{sec:pf}
% ===========================================================================

Since $R$ is proper, with poles only at $T=q^{-j}$, $0\le j\le n$, each of order at most $A$,
\begin{equation}\label{eq:pf}
R(T)=\sum_{j=0}^n\sum_{s=1}^A\frac{c_{s,j}(q)}{(1-q^jT)^s},\qquad c_{s,j}\in\QQ(q),
\end{equation}
and $\sum_{s,j}c_{s,j}=R(0)=0$, because $e\ge1$.

\begin{corollary}\label{cor:cinv}
$c_{s,j}(1/q)=q^{-e}c_{s,n-j}(q)$ for all $s,j$.
\end{corollary}

\begin{proof}
Apply the automorphism $q\mapsto1/q$ of $\QQ(q)$ to \eqref{eq:pf}: this gives the partial fractions of $R(T;1/q)$, with
poles at $T=q^j$. By Lemma~\ref{lem:inv} and \eqref{eq:pf}, $R(T;1/q)=q^{-e}\sum_{s,j'}c_{s,j'}(q)(1-q^{j'-n}T)^{-s}$.
Partial fractions are unique, and the pole $q^j$ corresponds to $j'=n-j$.
\end{proof}

Let $\Phi_\ell$ be the $\ell$-th cyclotomic polynomial and $d_n(q)=\prod_{\ell=1}^n\Phi_\ell(q)$. For $1\le m\le n$,
$1-q^m=-\prod_{\ell\mid m}\Phi_\ell(q)$ divides $d_n(q)$ in $\ZZ[q,q^{-1}]$, and so does $1-q^{-m}=-q^{-m}(1-q^m)$.

\begin{lemma}[denominators]\label{lem:den}
$d_n(q)^{A-s}\,c_{s,j}(q)\in\ZZ[q,q^{-1}]$ for all $1\le s\le A$ and $0\le j\le n$.
\end{lemma}

\begin{proof}
Since
\[
(q^{-rn}T;q)_{rn}=\prod_{\rho=1}^r(q^{-\rho n}T;q)_n\quad\text{and}\quad
(q^{n+1}T;q)_{rn}=\prod_{\rho=0}^{r-1}(q^{n+1+\rho n}T;q)_n,
\]
we can write $R(T)=T^e\prod_{i=1}^Af_i(T)$, where each $f_i$ is one of
\begin{itemize}
\item $(q;q)_n/(T;q)_{n+1}$ ($A-2r$ times);
\item $(q^hT;q)_n/(T;q)_{n+1}$ with $h=-\rho n$, $1\le\rho\le r$;
\item $(q^hT;q)_n/(T;q)_{n+1}$ with $h=n+1+\rho n$, $0\le\rho\le r-1$.
\end{itemize}
Each $f_i=N_i/(T;q)_{n+1}$ is proper with simple poles, so
\[
f_i(T)=\sum_{j'=0}^n\frac{a_{i,j'}}{1-q^{j'}T},\qquad a_{i,j'}=\frac{N_i(q^{-j'})}{\prod_{i'\ne j'}(1-q^{i'-j'})},
\]
where $i'$ runs over $0\le i'\le n$, and
\[
\prod_{i'\ne j'}(1-q^{i'-j'})=\pm q^{-j'(j'+1)/2}(q;q)_{j'}(q;q)_{n-j'}.
\]
Hence $a_{i,j'}=\pm q^{j'(j'+1)/2}\binom n{j'}_qN_i(q^{-j'})/(q;q)_n$, where $\binom n{j'}_q\in\ZZ[q]$ is the
Gaussian binomial coefficient. For the first type $N_i/(q;q)_n=1$. For the other two, $N_i(q^{-j'})=(q^b;q)_n$ with
$b=h-j'$; for the second type $b\le-n$, for the third $b\ge1$. If $b\ge1$, $(q^b;q)_n/(q;q)_n=\binom{b+n-1}n_q$; if
$b\le-n$, then $(q^b;q)_n=\pm q^{-N}(q^{c};q)_n$ with $c=1-b-n\ge1$ and some integer $N$, and the quotient is again
$\pm q^{-N}$ times a Gaussian binomial coefficient. So every $a_{i,j'}$ lies in $\ZZ[q,q^{-1}]$.

Fix $j$ and put $u=1-q^jT$, so $T=q^{-j}(1-u)$. Then $T^e=q^{-je}(1-u)^e$ has coefficients in $\ZZ[q^{\pm1}]$ as a
polynomial in $u$, and
\[
(1-q^jT)f_i(T)=a_{i,j}+u\sum_{j'\ne j}\frac{a_{i,j'}}{\alpha_{j'}+q^{j'-j}u},\qquad\alpha_{j'}=1-q^{j'-j}.
\]
The coefficient of $u^M$, $M\ge1$, in this power series is
\[
\sum_{j'\ne j}a_{i,j'}\,(-q^{j'-j})^{M-1}\,\alpha_{j'}^{-M}.
\]
Each $\alpha_{j'}$ is $1-q^{\pm m}$ with $1\le m\le n$, so $d_n(q)^M$ times this coefficient lies in $\ZZ[q^{\pm1}]$.
Finally $c_{s,j}$ is the coefficient of $u^{A-s}$ in $u^AR=T^e\prod_i(1-q^jT)f_i(T)$. Distributing the $u$-degree
$A-s$ over the factors, with degree $M_i$ on the factor $(1-q^jT)f_i$ and $\sum_iM_i\le A-s$, gives
$d_n^{A-s}c_{s,j}\in\ZZ[q^{\pm1}]$: the constant terms $a_{i,j}$ ($M_i=0$) and the coefficients of $T^e$ need no
denominator, and a factor with $M_i\ge1$ needs at most $d_n^{M_i}$.
\end{proof}

% ===========================================================================
\section{The linear forms}\label{sec:forms}
% ===========================================================================

From now on $q$ also denotes a complex number with $0<|q|<1$. The identities of Sections~\ref{sec:R}
and~\ref{sec:pf}, and those of Section~\ref{sec:recip} below, take place in the ring
$\QQ[q,q^{-1},d_n(q)^{-1}]$ (by Lemma~\ref{lem:den}, and since $1-q^{\pm m}$ divides $d_n(q)$ for $m\le n$), or in the
ring of rational functions of $T$ over it. Since every root of $d_n$ is a root of unity, evaluation at a number $q$ with
$q\ne0$ and $|q|\ne1$ is a ring homomorphism of $\QQ[q,q^{-1},d_n(q)^{-1}]$ to $\CC$. This ring is stable under
$q\mapsto1/q$, because $\Phi_1(1/q)=-q^{-1}\Phi_1(q)$ and $\Phi_\ell(1/q)=q^{-\eu(\ell)}\Phi_\ell(q)$ for $\ell\ge2$,
and the value at $q$ of $f(1/q)$ is the value of $f$ at $1/q$. So all these identities, including those relating
$f(1/q)$ and $f(q)$, may be
evaluated at such $q$. Define
\[
S_n(q)=\sum_{k\ge1}R(q^k)-\sum_{k\ge1}R(q^{-n-k}).
\]
No argument is a pole. By Lemma~\ref{lem:wp}, $R(q^{-n-k})=q^{A(n/2+k)}R(q^k)$, and $R(q^k)=O(|q|^{ke})$, so both
sums converge absolutely.

Put $C_s(x)=\binom{x+s-1}{s-1}\in\QQ[x]$ and $\mathcal Q(x)=\sum_{s,j}c_{s,j}C_s(x)=\sum_{\lambda}\mathcal
Q_\lambda x^\lambda$, a polynomial of degree at most $A-1$.

\begin{lemma}[expansion]\label{lem:expansion}
For $0<|q|<1$,
\[
S_n(q)=P_0(q)+\sum_{\substack{3\le l\le A-1\\ l\text{ odd}}}P_l(q)\,\zeta_q(l),
\]
where $P_l=2\mathcal Q_{l-1}$ and
\[
P_0=-\sum_{s,j}c_{s,j}W_{s,j},\qquad W_{s,j}=\sum_{m=1}^j\bigl((1-q^m)^{-s}-1\bigr)-\sum_{m=1}^{n-j}(1-q^{-m})^{-s}.
\]
\end{lemma}

\begin{proof}
\emph{Positive side.} Since $\sum_{s,j}c_{s,j}=0$, $R(q^k)=\sum_{s,j}c_{s,j}\bigl[(1-q^{j+k})^{-s}-1\bigr]$. Summing over
$k\ge1$ and using
\[
\sum_{m\ge1}\bigl[(1-q^m)^{-s}-1\bigr]=\sum_{m\ge1}\sum_{i\ge1}C_s(i)q^{mi}=\sum_{i\ge1}C_s(i)\frac{q^i}{1-q^i},
\]
we get $\sum_{k\ge1}R(q^k)=\sum_{i\ge1}\mathcal Q(i)\frac{q^i}{1-q^i}-\sum_{s,j}c_{s,j}\sum_{m=1}^j\bigl((1-q^m)^{-s}
-1\bigr)$.

\emph{Negative side.} We have
\[
\sum_{k\ge1}(1-q^{j-n-k})^{-s}=\sum_{m\ge1}(1-q^{-m})^{-s}-\sum_{m=1}^{n-j}(1-q^{-m})^{-s}.
\]
Moreover $(1-q^{-m})^{-s}=(-1)^sq^{ms}(1-q^m)^{-s}=(-1)^s\sum_{i\ge s}\binom{i-1}{s-1}q^{mi}$, and
$(-1)^s\binom{i-1}{s-1}=-C_s(-i)$; both sides vanish for $1\le i<s$. Hence
$\sum_{m\ge1}(1-q^{-m})^{-s}=-\sum_{i\ge1}C_s(-i)\frac{q^i}{1-q^i}$ and
\[
\sum_{k\ge1}R(q^{-n-k})=-\sum_{i\ge1}\mathcal Q(-i)\frac{q^i}{1-q^i}-\sum_{s,j}c_{s,j}\sum_{m=1}^{n-j}(1-q^{-m})^{-s}.
\]

\emph{Difference.} The $q$-zeta part of $S_n$ is
\[
\sum_{i\ge1}\bigl[\mathcal Q(i)+\mathcal Q(-i)\bigr]\frac{q^i}{1-q^i}=\sum_{\lambda\text{ even}}2\mathcal
Q_\lambda\,\zeta_q(\lambda+1).
\]
Here $\mathcal Q_0=\mathcal Q(0)=\sum_{s,j}c_{s,j}=0$,
and $\lambda\le A-2$ since $\lambda$ is even and at most $A-1$. All rearrangements are of absolutely convergent double
series.
\end{proof}

The values $\zeta_q(l)$ of even weight are absent from $S_n$: the minus sign between the two sums keeps only the even
part of $\mathcal Q$.

\begin{remark}
Not used below. Put $S_n^+(q)=\sum_{k\ge1}R(q^k;q)$. By Lemma~\ref{lem:inv}, $R(q^{-n-i};q)=q^eR(q^{-i};1/q)$, so
$S_n=S_n^+(q)-q^eS_n^+(1/q)$, where $S_n^+(1/q)$ means $\sum_{i\ge1}R(q^{-i};1/q)$. The combination $S_n^{[1]}$ of
\cite[\S3]{KRZ06} has the same shape, but is built from $TR(T)$ in place of $R(T)$; it is not a constant multiple of
our $S_n$.
\end{remark}

% ===========================================================================
\section{Reciprocity}\label{sec:recip}
% ===========================================================================

\begin{lemma}\label{lem:recip}
As elements of $\QQ(q)$: $P_l(1/q)=q^{-e}P_l(q)$ for odd $l\ge3$, but $P_0(1/q)=-q^{-e}P_0(q)$.
\end{lemma}

\begin{proof}
By Corollary~\ref{cor:cinv},
\[
\mathcal Q(x;1/q)=\sum_{s,j}c_{s,j}(1/q)C_s(x)=q^{-e}\sum_{s,j}c_{s,n-j}(q)C_s(x)=q^{-e}\mathcal Q(x;q),
\]
which gives the first claim. For $P_0$: $W_{s,n-j}$ evaluated at $1/q$ is
\[
\sum_{m=1}^{n-j}\bigl((1-q^{-m})^{-s}-1\bigr)-\sum_{m=1}^j(1-q^m)^{-s}=-W_{s,j}(q)-n.
\]
Hence, reindexing $j\mapsto n-j$ and using Corollary~\ref{cor:cinv},
\begin{align*}
P_0(1/q)&=-\sum_{s,j}c_{s,n-j}(1/q)\,W_{s,n-j}(1/q)\\
&=-q^{-e}\sum_{s,j}c_{s,j}(q)\bigl(-W_{s,j}(q)-n\bigr)=-q^{-e}P_0(q),
\end{align*}
using $\sum_{s,j}c_{s,j}=0$ once more.
\end{proof}

\begin{remark}
Heuristically, $q\mapsto1/q$ sends $\zeta_q(l)$ to $-\zeta_q(l)-\zeta(1-l)$, and $\zeta(1-l)=0$ for odd $l\ge3$: so
$S_n$ behaves like an eigenvector of $q\mapsto1/q$ if the $q$-zeta values are treated as anti-invariant. At a unit with
$\sigma q=1/q$, Galois conjugation realises this for the numbers $X$ with $\sigma X=-X$, that is, for $X\in\QQ\sqrt d$.
This remark is not used in the proofs.
\end{remark}

% ===========================================================================
\section{Size and nonvanishing}\label{sec:size}
% ===========================================================================

\begin{lemma}\label{lem:size}
Fix $q$ with $0<|q|<1$. For all sufficiently large $n$, $S_n(q)\ne0$ and
\[
\log|S_n(q)|=-(rn+1)e\log(1/|q|)+O_q(1);
\]
in particular $\log|S_n(q)|=-\frac12r(A-2r)\,n^2\log(1/|q|)+O_q(n)$.
\end{lemma}

\begin{proof}
Let $\kappa=|q|$, $c_-=\prod_{i\ge1}(1-\kappa^i)>0$ and $c_+=\prod_{i\ge1}(1+\kappa^i)$. For $a\ge1$ and $m\ge0$,
$c_-\le|(q^a;q)_m|\le c_+$. For $1\le k\le rn$, $R(q^k)=0$, because $(q^{k-rn};q)_{rn}$ contains the factor $1-q^0$.
For $k\ge rn+1$ every $q$-Pochhammer symbol in $R(q^k)$ has base $q^a$ with $a\ge1$; there are $A-2r+2\le A$ of them in
the numerator and $A$ in the denominator, so $C^{-1}\kappa^{ke}\le|R(q^k)|\le C\kappa^{ke}$ with $C=(c_+/c_-)^{2A}$,
a constant depending only on $\kappa$ and $A$.
By Lemma~\ref{lem:wp},
\[
S_n=\sum_{k\ge rn+1}\bigl(1-q^{A(n/2+k)}\bigr)R(q^k),
\]
and $|1-q^{A(n/2+k)}|\in[\frac12,2]$ for these $k$ once $n$ is large. So
\[
\bigl|S_n-(1-q^{A(n/2+rn+1)})R(q^{rn+1})\bigr|\le\frac{2C\kappa^{(rn+2)e}}{1-\kappa^e},
\]
while
\[
\bigl|(1-q^{A(n/2+rn+1)})R(q^{rn+1})\bigr|\ge\frac{\kappa^{(rn+1)e}}{2C}.
\]
The ratio of the two bounds is $4C^2\kappa^e/(1-\kappa^e)\to0$, since $e\to\infty$. Hence $|S_n|$ lies between
$\kappa^{(rn+1)e}/(4C)$ and $4C\kappa^{(rn+1)e}$ for large $n$.
\end{proof}

Krattenthaler, Rivoal and Zudilin obtain the analogous limit for their series by a saddle-point argument
\cite{KRZ06}; we do not need it.

% ===========================================================================
\section{Cyclotomic norms at a unit}\label{sec:cost}
% ===========================================================================

From now on $q$ is as in Theorem~\ref{thm:main}, and $\kappa=|q|$. Since $d_n$ has integer coefficients and
$\sigma q=1/q$, the norm of $d_n(q)$ is $\Nm d_n(q)=d_n(q)\,d_n(1/q)$. Write $\eu$ for Euler's function (the letter $\varphi$ denotes the golden ratio) and
$\Tot(n)=\sum_{\ell\le n}\eu(\ell)$.

\begin{lemma}\label{lem:cost}
With $B=\sum_{a\ge1}-\log(1-\kappa^a)<\infty$,
\[
\bigl|\log|\Nm d_n(q)|-\Tot(n)\log(1/\kappa)\bigr|\le2Bn,\qquad\text{and}\qquad\Tot(n)\le\tfrac38(n+1)^2.
\]
\end{lemma}

\begin{proof}
From $\Phi_\ell(x)=\prod_{a\mid\ell}(x^a-1)^{\mu(\ell/a)}$ and $\sum_{a\mid\ell}a\,\mu(\ell/a)=\eu(\ell)$,
\[
|\Phi_\ell(q)\Phi_\ell(1/q)|=\kappa^{-\eu(\ell)}\prod_{a\mid\ell}|1-q^a|^{2\mu(\ell/a)}.
\]
Since $1-\kappa^a\le|1-q^a|\le1+\kappa^a$, we have $|\log|1-q^a||\le-\log(1-\kappa^a)$, so
$|\log|\Nm\Phi_\ell(q)|-\eu(\ell)\log(1/\kappa)|\le2B$. Sum over $\ell\le n$.

For the second claim, $\eu(\ell)\le\ell$ for odd $\ell$, and $\eu(2m)\le m$ (if $m$ is odd, $\eu(2m)=\eu(m)\le m$; if
$m$ is even, $\eu(2m)=2\eu(m)\le m$). Hence
\[
\Tot(n)\le\sum_{\substack{\ell\le n\\ \ell\text{ odd}}}\ell+\sum_{m\le n/2}m\le\frac{(n+1)^2}4+\frac{n(n+2)}8
=\frac{3n^2+6n+2}8\le\frac38(n+1)^2.\qedhere
\]
\end{proof}

\begin{remark}\label{rem:mertens}
By the classical asymptotic $\Tot(n)=3n^2/\pi^2+O(n\log n)$ \cite[Theorem~330]{HW08}, Lemma~\ref{lem:cost} gives
$\log|\Nm d_n(q)|=(3/\pi^2)n^2\log(1/\kappa)+O(n\log n)$. The proof below only needs $\Tot(n)\le cn^2+O(n)$ for some
$c<\frac12$; the trivial bound $\Tot(n)\le n(n+1)/2$ would not suffice.
\end{remark}

% ===========================================================================
\section{Proof of the main results}\label{sec:proof}
% ===========================================================================

\begin{proof}[Proof of Theorem~\ref{thm:main}]
Take $A=4$ and $r=1$, so $e=n$. By Lemma~\ref{lem:expansion}, $S_n=P_0+P_3\zeta_q(3)$, and since $[x^2]C_3=\frac12$,
$[x^2]C_4=1$ and $C_1,C_2$ have degree less than $2$,
\[
P_3=2\mathcal Q_2=\sum_{j=0}^n(c_{3,j}+2c_{4,j}).
\]
By Lemma~\ref{lem:den}, $d_n(q)P_3\in\ZZ[q^{\pm1}]$. Also $d_n(q)^sW_{s,j}\in\ZZ[q^{\pm1}]$, because $1-q^{\pm m}$
divides $d_n(q)$ for $m\le n$; with Lemma~\ref{lem:den} this gives $d_n(q)^4P_0\in\ZZ[q^{\pm1}]$. Put
$\delta_n=d_n(q)^4$. Since $q$ is a unit, $\ZZ[q^{\pm1}]\subset\cO_K$, so $\delta_nP_0$ and $\delta_nP_3$ are
values at $q$ of Laurent polynomials with integer coefficients, and lie in $\cO_K$.

Let $\theta\in\cO_K$ be nonzero with $\sigma\theta=-\theta$; then $\theta^2$ is a nonzero rational integer. We show
that $\theta\zeta_q(3)\notin\QQ$; the theorem follows with $\theta=\sqrt d$. Suppose $\theta\zeta_q(3)=a/b$ with
$a,b\in\ZZ$, $b\ge1$. Then $\zeta_q(3)=a\theta/(b\theta^2)$ and
\[
y_n:=b\theta^2\,\delta_nS_n=b\theta^2\,\delta_nP_0+a\theta\,\delta_nP_3\in\cO_K.
\]
For large $n$, $y_n\ne0$ by Lemma~\ref{lem:size}, since $d_n(q)\ne0$. Apply $\sigma$. For a Laurent polynomial $f$ with
integer coefficients, $\sigma(f(q))=f(1/q)$; with Lemma~\ref{lem:recip}, evaluated at $q$ as explained in
Section~\ref{sec:forms}, and $\sigma\theta=-\theta$,
\[
\sigma(y_n)=d_n(1/q)^4\bigl[b\theta^2\cdot(-q^{-n}P_0)-a\theta\cdot q^{-n}P_3\bigr]=-q^{-n}\,d_n(1/q)^4\,b\theta^2\,S_n.
\]
Hence
\[
|\Nm y_n|=|y_n\,\sigma(y_n)|=b^2\theta^4\,\kappa^{-n}\,|\Nm d_n(q)|^4\,|S_n|^2.
\]
By Lemmas~\ref{lem:size} and~\ref{lem:cost},
\[
\log|\Nm y_n|\le\Bigl(4\cdot\frac38-2\Bigr)n^2\log(1/\kappa)+O_q(n)=-\frac12\,n^2\log(1/\kappa)+O_q(n),
\]
which is negative for large $n$. This contradicts $0\ne y_n\in\cO_K$, for which $|\Nm y_n|\ge1$.
\end{proof}

\begin{remark}
With Remark~\ref{rem:mertens} in place of the elementary bound, the exponent is $12/\pi^2-2=-0.7841\ldots$. Numerically,
with the exact integer $\Nm d_n(q)$ and $S_n$ summed directly, the quantity
$(\log|\Nm\delta_n|+2\log|S_n|+n\log(1/\kappa))/(n^2\log(1/\kappa))$ stays near this value for
$6\le n\le30$: it lies in $[-0.82,-0.58]$ at $q=\varphi^{-2}$ and in $[-0.84,-0.72]$ at $q=-\varphi^{-2}$. This is not
used in the proof. An irrationality measure for $\sqrt d\,\zeta_q(3)$ would in addition need an upper bound for
$|P_3(q)|$, as in Lemma~\ref{lem:coef}; we do not pursue it.
\end{remark}

\begin{proof}[Proof of Corollary~\ref{cor:lucas}]
Let $(P,Q)$ be admissible, $D=P^2-4Q$, $\eta=(P+\sqrt D)/2$ and $\bar\eta=(P-\sqrt D)/2=Q/\eta$. Then $\eta\ge\varphi>1$
(for $Q=-1$, $\eta\ge(1+\sqrt5)/2$; for $Q=1$, $\eta\ge(3+\sqrt5)/2$) and $|\bar\eta|=1/\eta<1$, and $U_m=(\eta^m-\bar\eta^m)/\sqrt D$ by Binet's formula, so
$U_m\ne0$ for $m\ge1$. Moreover
\[
U_{m+1}-\eta U_m=\frac{\bar\eta^m(\eta-\bar\eta)}{\sqrt D}=\bar\eta^m.
\]
Put $q=\bar\eta/\eta=Q\eta^{-2}$. Then $0<|q|<1$, $q$ is a unit of $\ZZ[\eta]$, and $\sigma q=\eta/\bar\eta=1/q$, where
$\sigma$ is the nontrivial automorphism of $K=\QQ(\sqrt D)$; so $q$ is a unit of norm $+1$ of $K$. Since
$q^m/(1-q^m)=\bar\eta^m/(\eta^m-\bar\eta^m)=\bar\eta^m/(\sqrt D\,U_m)$,
\[
\sum_{m\ge1}m^2\Bigl(\frac{U_{m+1}}{U_m}-\eta\Bigr)=\sum_{m\ge1}\frac{m^2\bar\eta^m}{U_m}=\sqrt D\,\zeta_q(3),
\]
which is irrational by Theorem~\ref{thm:main} with $d=D$.
\end{proof}

\begin{proof}[Proof of Corollary~\ref{cor:fib}]
Take $(P,Q)=(1,-1)$, where $\eta=\varphi$, $\bar\eta=-\varphi^{-1}$ and $q=-\varphi^{-2}$, and $(P,Q)=(2,-1)$, where
$\eta=1+\sqrt2$. The second expression for the Fibonacci sum is $\sum_mm^2\bar\eta^m/F_m$.
\end{proof}

\begin{remark}\label{rem:plain}
The method does not show that $\zeta_q(3)\notin\QQ$. If $\zeta_q(3)$ were rational, then by Lemma~\ref{lem:recip}
$\sigma(S_n)=q^{-n}(S_n-2P_0)$, which is not small: it has the size of $P_0$. Numerically $|P_0|$ is about
$|q|^{-n^2}$ (we find $\log|P_0|/(n^2\log(1/|q|))\in[0.90,0.95]$ at $n=12$ for the units $q=\pm\varphi^{-2}$, $\pm(2-\sqrt3)$, $\pm(3-2\sqrt2)$ and $(5-\sqrt{21})/2$), which is the size of
the upper bound of Lemma~\ref{lem:coef} at $A=4$, $r=1$ (a matching lower bound, for the series of \cite{KRZ06}, is the
conjecture of \cite[\S3.5]{KRZ06}); with it the norm estimate fails. Heuristically this is the obstruction met at $q=1/p$, where
clearing the $p$-power denominators of the coefficients costs as much.
The concurrent preprint \cite{Twi26} claims $\zeta_q(3)\notin\QQ(q)$, which contains this plain statement
(see the introduction).
\end{remark}

\begin{remark}\label{rem:fibzeta}
Theorem~\ref{thm:main} and the irrationality of $\zeta_F(3)=\sum_mF_m^{-3}$ \cite{Kri26a} do not imply each other, as far
as we can see. Expanding $F_m^{-3}$ in powers of $\varphi^{-m}$ and splitting the resulting double series by the parity
of the outer index expresses $\zeta_F(3)$ as a $\QQ(\sqrt5)$-linear combination of the numbers
$\zeta_q(l)-2^{l-1}\zeta_{q^2}(l)$, $l\in\{1,3\}$, at $q=\varphi^{-2}$, and of values of Eisenstein series of odd weight
with the character $\chi_{-4}$ (cf.\ \cite[\S1]{Kri26a}). The linear forms used here see only the $q$-zeta part, so the
method gives no information on $\zeta_F(3)$.
\end{remark}

% ===========================================================================
\section{Many values}\label{sec:many}
% ===========================================================================

% Source: handoff/fibzeta/many/PROOF_T2.md (Lemmas 8--11, sections 5--7). Not Lean-checked.
In this section $A\ge4$ is even and $1\le r<A/2$ are arbitrary, $q$ is as in Theorem~\ref{thm:main}, $\kappa=|q|$,
$\nu=\log(1/\kappa)>0$ and $e=(A-2r)n/2$. We prove Theorem~\ref{thm:many}. Besides the lemmas of
Sections~\ref{sec:R}--\ref{sec:cost}, which hold for all such $(A,r)$, it uses two classical inputs: the asymptotic
$\Tot(n)=3n^2/\pi^2+O(n\log n)$ of Remark~\ref{rem:mertens}, in both directions, and Nesterenko's linear independence
criterion. None of this section is formalised.

\subsection*{Denominators}

\begin{lemma}\label{lem:dengen}
In $\QQ(q)$, $(A-1)!\,d_n(q)^{A-l}P_l\in\ZZ[q^{\pm1}]$ for odd $3\le l\le A-1$, and $d_n(q)^AP_0\in\ZZ[q^{\pm1}]$.
Hence $\Delta_n(q):=(A-1)!\,d_n(q)^A$ satisfies $\Delta_nP_0,\Delta_nP_l\in\ZZ[q^{\pm1}]$.
\end{lemma}

\begin{proof}
Since $(s-1)!\,C_s(x)=(x+1)(x+2)\cdots(x+s-1)\in\ZZ[x]$, $(A-1)!\,[x^\lambda]C_s\in\ZZ$ for $s\le A$, and
$[x^\lambda]C_s=0$ unless $\lambda\le s-1$. So only $s\ge l$ contributes to $P_l=2\sum_{s,j}c_{s,j}[x^{l-1}]C_s$, and
for those $s$,
\[
(A-1)!\,d_n^{A-l}c_{s,j}\cdot2[x^{l-1}]C_s=d_n^{s-l}\cdot\bigl(d_n^{A-s}c_{s,j}\bigr)\cdot2(A-1)!\,[x^{l-1}]C_s
\in\ZZ[q^{\pm1}]
\]
by Lemma~\ref{lem:den}. For $P_0$, $d_n^sW_{s,j}\in\ZZ[q^{\pm1}]$ because $1-q^{\pm m}$ divides $d_n$ for $m\le n$, and
$d_n^Ac_{s,j}W_{s,j}=(d_n^{A-s}c_{s,j})(d_n^sW_{s,j})$.
\end{proof}

As in the proof of Theorem~\ref{thm:main}, the values $\Delta_nP_0$ and $\Delta_nP_l$ at $q$ lie in $\ZZ[q^{\pm1}]\subset
\cO_K$, and by Lemma~\ref{lem:recip}
\begin{equation}\label{eq:sigmaP}
\sigma\bigl(P_l(q)\bigr)=q^{-e}P_l(q),\qquad\sigma\bigl(P_0(q)\bigr)=-q^{-e}P_0(q),
\end{equation}
since $\sigma$ maps the value at $q$ of a quotient of integer Laurent polynomials to its value at $1/q$, and the
denominators $\Delta_n(q)$, $\Delta_n(1/q)$ do not vanish ($q$ is not a root of unity).

\subsection*{Coefficients}

\begin{lemma}\label{lem:coef}
Let $q\in\CC$ with $0<\kappa=|q|<1$ and $\rho=(1-\kappa)/2$. Put $c_+=\prod_{i\ge1}(1+\kappa^i)$,
$\gamma_-=\prod_{i\ge1}(1-(1+\rho)\kappa^i)$ and $\gamma_+=\prod_{i\ge1}(1+(1+\rho)\kappa^i)$. Then $\gamma_->0$, and for
$0\le j\le n$, $1\le s\le A$,
\[
|c_{s,j}(q)|\le\rho^{s-A(j+1)}\,c_+^{A-2r+1}\,\gamma_+\,\gamma_-^{-A}\,(1+\rho)^{rn+e}\,\kappa^{-\mathcal E(j)},
\]
where
\[
\mathcal E(j)=\tfrac A2j(n-j)-\tfrac A2j+\tfrac{r^2n^2}2+\tfrac{rn}2.
\]
Consequently $\log\max\bigl(|P_0(q)|,\max_l|P_l(q)|\bigr)\le\frac{A+4r^2}8\,n^2\log(1/\kappa)+O(n)$, with the
implied constant depending only on $A$, $r$ and $\kappa$.
\end{lemma}

\begin{proof}
Put $u=1-q^jT$ and $w=1-u$, so $T=q^{-j}w$. As a function of $u$, $R$ has its poles at $u_i=1-q^{j-i}$, $0\le i\le n$,
and $|u_i|\ge1-\kappa>\rho$ for $i\ne j$. By \eqref{eq:pf} the principal part of $R$ at $u=0$ is
$\sum_sc_{s,j}u^{-s}$, so $c_{s,j}=\frac1{2\pi i}\oint_{|u|=\rho}R\,u^{s-1}\,du$ and $|c_{s,j}|\le\rho^s\max_{|u|=\rho}
|R|$. On this circle $1-\rho\le|w|\le1+\rho$; note $1-\rho-\kappa=\rho$ and $(1+\rho)\kappa=\kappa(3-\kappa)/2<1$, which
gives $\gamma_->0$. Bound the factors of \eqref{eq:R}:
\begin{itemize}
\item $|(q;q)_n|^{A-2r}\le c_+^{A-2r}$;
\item $(q^{-rn}T;q)_{rn}=\prod_{M=j+1}^{j+rn}(1-q^{-M}w)$, and $|1-q^{-M}w|\le\kappa^{-M}(1+\rho)(1+\kappa^M)$, so
the modulus is at most $\kappa^{-rnj-rn(rn+1)/2}(1+\rho)^{rn}c_+$;
\item $(q^{n+1}T;q)_{rn}=\prod_{i=0}^{rn-1}(1-q^{n+1+i-j}w)$ has distinct exponents $\ge1$, so modulus at most $\gamma_+$;
\item $|T^e|\le\kappa^{-je}(1+\rho)^e$;
\item $(T;q)_{n+1}=\prod_{i=0}^n(1-q^{i-j}w)$: for $i<j$, $|1-q^{i-j}w|\ge\kappa^{i-j}(|w|-\kappa^{j-i})\ge
\kappa^{i-j}\rho$; for $i=j$ the factor is $u$, of modulus $\rho$; the factors with $i>j$ have product at least $\gamma_-$.
So $|(T;q)_{n+1}|\ge\kappa^{-j(j+1)/2}\rho^{j+1}\gamma_-$.
\end{itemize}
Multiplying, $|R|\le c_+^{A-2r+1}\gamma_+(1+\rho)^{rn+e}\kappa^{-\mathcal E(j)}\rho^{-A(j+1)}\gamma_-^{-A}$, where
$\mathcal E(j)=rnj+rn(rn+1)/2+je-Aj(j+1)/2$; since $rnj+je=Anj/2$, this is the stated $\mathcal E(j)$, and
$\mathcal E(j)\le An^2/8+r^2n^2/2+rn/2$. Since $\rho^{s-A(j+1)}\le\rho^{-A(n+1)}$, this gives
$\log\max_{s,j}|c_{s,j}|\le\frac{A+4r^2}8n^2\nu+O(n)$. The coefficients of $C_s$ are nonnegative with sum $C_s(1)=s$, so
$|P_l|\le(n+1)A(A+1)\max|c_{s,j}|$. For $1\le m\le n$, $|(1-q^m)^{-s}-1|\le(1-\kappa)^{-A}+1$ and
$|(1-q^{-m})^{-s}|\le\max(1,(\kappa/(1-\kappa))^A)$, so $|W_{s,j}|\le nw_0$ with $w_0$ depending on $A$ and $\kappa$
only, and $|P_0|\le n(n+1)Aw_0\max|c_{s,j}|$.
\end{proof}

This is the analogue for our coefficients of \cite[Lemma~5]{KRZ06}, proved by the same Cauchy estimate; we include the proof
because our coefficients $c_{s,j}$ are not those of \cite{KRZ06}, and because we want an explicit bound.

\subsection*{Balanced trace forms}

Since $\Phi_1(1/q)=-q^{-1}\Phi_1(q)$ and $\Phi_\ell(1/q)=q^{-\eu(\ell)}\Phi_\ell(q)$ for $\ell\ge2$ ($\Phi_\ell$ is
self-reciprocal of degree $\eu(\ell)$), we have exactly $\sigma(d_n(q))=d_n(1/q)=-q^{-\Tot(n)}d_n(q)$, so
$|\sigma\Delta_n(q)|=\kappa^{-A\Tot(n)}|\Delta_n(q)|$. Moreover $|\log|d_n(q)||\le\sum_{\ell\le n}\sum_{a\mid\ell}
|\log|1-q^a||\le nB$, with $B$ as in Lemma~\ref{lem:cost}, because each $a$ divides at most $n$ of the $\ell\le n$.
With Remark~\ref{rem:mertens},
\begin{equation}\label{eq:twoplaces}
\log|\Delta_n(q)|=O(n),\qquad\log|\sigma\Delta_n(q)|=\frac{3A}{\pi^2}n^2\nu+O(n\log n).
\end{equation}

In this subsection $d$ is squarefree; the general case is reduced to it at the end of the proof of Theorem~\ref{thm:many}. The point of the construction is to split the cost $\log|\Nm\Delta_n|$ evenly between the two real places. For
$n$ put $\varrho_n=\kappa^{-e}|\sigma\Delta_n|/|\Delta_n|$, and let $m=m_n\in\ZZ$ be the unique integer with
\begin{equation}\label{eq:balance}
2\le\kappa^{2m}\varrho_n<2\kappa^{-2}.
\end{equation}
Put $\Delta_n'=q^{-m}\Delta_n(q)\in\cO_K$. Since $\sigma q=1/q$, $|\Delta_n'|=\kappa^{-m}|\Delta_n|$ and
$|\sigma\Delta_n'|=\kappa^m|\sigma\Delta_n|$, so $M_n:=\kappa^{-e}|\sigma\Delta_n'|$ satisfies $|\Delta_n'|\le M_n/2$
and $M_n^2=\kappa^{2m}\varrho_n\cdot\kappa^{-e}|\Delta_n\sigma\Delta_n|\in\kappa^{-e}|\Delta_n\sigma\Delta_n|\cdot
[2,2\kappa^{-2})$. By \eqref{eq:twoplaces},
\begin{equation}\label{eq:Mn}
\log M_n=\frac{3A}{2\pi^2}n^2\nu+O(n\log n).
\end{equation}
For odd $3\le l\le A-1$ define
\[
\ell_{0,n}=\operatorname{Tr}_{K/\QQ}\bigl(\Delta_n'P_0(q)\bigr),\qquad
\ell_{l,n}=\operatorname{Tr}_{K/\QQ}\bigl(\Delta_n'P_l(q)/\sqrt d\bigr),
\]
\[
L_n=\ell_{0,n}+\sum_l\ell_{l,n}\,\sqrt d\,\zeta_q(l),
\]
where $K=\QQ(\sqrt d)$ with $d$ squarefree.

\begin{lemma}\label{lem:trace}
\textup{(a)} $\ell_{0,n},\ell_{l,n}\in\ZZ$.
\textup{(b)} $L_n=(\Delta_n'-q^{-e}\sigma\Delta_n')\,S_n(q)$.
\textup{(c)} For large $n$, $L_n\ne0$ and $\log|L_n|=\bigl(\frac{3A}{2\pi^2}-\frac{r(A-2r)}2\bigr)n^2\nu+O(n\log n)$.
\textup{(d)} $\log\max(|\ell_{0,n}|,\max_l|\ell_{l,n}|)\le\bigl(\frac{3A}{2\pi^2}+\frac{A+4r^2}8\bigr)n^2\nu+
O(n\log n)$.
\end{lemma}

\begin{proof}
(a) $\operatorname{Tr}(\cO_K)\subset\ZZ$, and $\operatorname{Tr}(x/\sqrt d)\in\ZZ$ for $x\in\cO_K$: if
$d\equiv2,3\pmod4$ then $x=a+b\sqrt d$ with $a,b\in\ZZ$ and the trace is $2b$; if $d\equiv1\pmod4$ then
$x=(a+b\sqrt d)/2$ with $a\equiv b\pmod 2$ and the trace is $b$. Apply this to $x=\Delta_n'P_0(q)$ and
$x=\Delta_n'P_l(q)$, which lie in $\cO_K$ by Lemma~\ref{lem:dengen}.

(b) Write $\operatorname{Tr}(z)=z+\sigma z$. The terms $z$ add up to $\Delta_n'(P_0+\sum_lP_l\zeta_q(l))=
\Delta_n'S_n$ by Lemma~\ref{lem:expansion}. By \eqref{eq:sigmaP} and $\sigma\sqrt d=-\sqrt d$, the conjugate terms are
$\sigma(\Delta_n'P_0)=-q^{-e}\sigma(\Delta_n')P_0$ and $\sigma(\Delta_n'P_l/\sqrt d)\cdot\sqrt d\,\zeta_q(l)=
-q^{-e}\sigma(\Delta_n')P_l\zeta_q(l)$; they add up to $-q^{-e}\sigma(\Delta_n')S_n$.

(c) $|\Delta_n'-q^{-e}\sigma\Delta_n'|\in[M_n-|\Delta_n'|,M_n+|\Delta_n'|]\subset[M_n/2,3M_n/2]$. Combine with
Lemma~\ref{lem:size} and \eqref{eq:Mn}.

(d) $|\ell_{0,n}|\le(|\Delta_n'|+M_n)|P_0(q)|\le\frac32M_n|P_0(q)|$, and likewise $|\ell_{l,n}|\le\frac32M_n|P_l(q)|/
\sqrt d$. Apply Lemma~\ref{lem:coef} and \eqref{eq:Mn}.
\end{proof}

\subsection*{Proof of Theorem~\ref{thm:many}}
We use Nesterenko's criterion \cite{Nes85}, in the precise form of \cite[Theorem~A]{FZ10} (under its
hypothesis~(N)), rescaled as in \cite[Proposition~2]{KRZ06}: let $\vartheta_1,\dots,
\vartheta_N$ be real numbers, and suppose that there are linear forms $\sum_jp_{j,n}\vartheta_j$ with
$p_{j,n}\in\ZZ$ such that, as $n\to\infty$, $\frac1{n^2}\log|\sum_jp_{j,n}\vartheta_j|\to\alpha<0$ and
$\limsup\frac1{n^2}\log|p_{j,n}|\le\beta$ for every $j$. Then $\dim_\QQ(\QQ\vartheta_1+\dots+\QQ\vartheta_N)\ge
1-\alpha/\beta$. (The usual statement has $N'$ in place of $n^2$; it applies to the forms indexed by
$N'\mapsto\lfloor\sqrt{N'}\rfloor$, since $\lfloor\sqrt{N'}\rfloor^2=N'+O(\sqrt{N'})$, and it does not require the
forms to be distinct.)

Take $\vartheta_1=1$ and $\vartheta_2,\dots,\vartheta_{A/2}$ equal to $\sqrt d\,\zeta_q(3),\dots,\sqrt d\,
\zeta_q(A-1)$, with the forms $L_n$. By Lemma~\ref{lem:trace},
$\alpha=(\frac{3A}{2\pi^2}-\frac{r(A-2r)}2)\nu$ and $\beta=(\frac{3A}{2\pi^2}+\frac{A+4r^2}8)\nu>0$. If $\alpha\ge0$
then $\tau(A,r)\le0$ and there is nothing to prove; otherwise
\[
\dim\ge1-\frac\alpha\beta=1+\frac{r(A-2r)/2-3A/(2\pi^2)}{3A/(2\pi^2)+(A+4r^2)/8}=1+\tau(A,r).
\]
The bound does not depend on $q$ or $d$, and it is the same for every $d$ with $K=\QQ(\sqrt d)$, since
$\sqrt d\,\zeta_q(l)$ changes by a common rational factor. \qed

\begin{proof}[Proof of the asymptotic bound and of Corollary~\ref{cor:many}]
Put $r=t\sqrt A$. Then $\tau(A,t\sqrt A)/\sqrt A\to4\pi^2t/(12+\pi^2+4\pi^2t^2)$ as $A\to\infty$, which is largest at
$t=\sqrt{12+\pi^2}/(2\pi)=0.744\ldots$, with value $\pi/\sqrt{12+\pi^2}=0.67178\ldots$; taking
$r=\lceil t\sqrt A\rceil$, which is admissible for $A\ge6$, gives the same limit. For Corollary~\ref{cor:many},
$1+\tau(12,2)=1+(64\pi^2-144)/(144+28\pi^2)=2.1601\ldots>2$, so
$\dim_\QQ(\QQ+\QQ\sqrt d\zeta_q(3)+\dots+\QQ\sqrt d\zeta_q(11))\ge3$, while $\dim_\QQ(\QQ+\QQ\sqrt d\zeta_q(3))\le2$.
The Lucas form follows from the identity $\sum_mm^{l-1}(U_{m+1}/U_m-\eta)=\sqrt D\,\zeta_q(l)$, proved for $l=3$ in the
proof of Corollary~\ref{cor:lucas} and valid for every $s$ by the same computation.
\end{proof}

\begin{remark}\label{rem:balance}
Without the balancing, that is with $\Delta_n$ in place of $\Delta_n'$, the form carries the whole cost
$\log|\Nm\Delta_n|$ at one place, and the same argument gives only
$1+(4\pi^2r(A-2r)-24A)/(24A+\pi^2(A+4r^2))$, which is asymptotically $(\pi/\sqrt{24+\pi^2})\sqrt A=0.5398\ldots\sqrt A$
and does not recover Theorem~\ref{thm:main} at $A=4$. With the balancing, writing $\mathrm{dec}=r(A-2r)/2$,
$\mathrm{deg}=3A/\pi^2$ and $\mathrm{coef}=(A+4r^2)/8$, we have $1+\tau=(\mathrm{dec}+\mathrm{coef})/(\mathrm{deg}/2+
\mathrm{coef})$, while the bound of \cite[Theorem~2]{KRZ06} at $q=1/p$ is $(\mathrm{dec}+\mathrm{coef})/(\mathrm{deg}+
2\,\mathrm{coef})$, where the second place is $p$-adic and the denominator must also clear powers of $p$ of size
$\mathrm{coef}$. So $1+\tau(A,r)$ is exactly twice the bound of \cite{KRZ06}. At $A=4$, $r=1$ the condition $\tau>0$
reads $2>12/\pi^2$, the inequality behind the proof of Theorem~\ref{thm:main} with the asymptotic of
Remark~\ref{rem:mertens}.

The balancing \eqref{eq:balance} itself is the standard device of moving an $\cO_K$-multiplier by a unit so that its
archimedean sizes are equal, as in Siegel's lemma over a number field \cite[Chapter~2]{BG06}. What is specific here is
its combination with the Galois-equivariance \eqref{eq:sigmaP} and the twist by $\sqrt d$, which make the form small at
both real places.

The comparison with \cite{KRZ06} is one of the same leading constants (decay, coefficient size and denominator, which our variant of the series shares with that of \cite{KRZ06}) and the same single-form criterion, at two kinds of $q$. At $q=1/p$ the results of
\cite{KRZ06} have been improved: Jouhet and Mosaki \cite{JM10} prove that one power of $d_n$ can be saved in the
denominators of a variant $\tilde S_n$ of the series, and deduce that one of $\zeta_q(3),\zeta_q(5),\zeta_q(7),
\zeta_q(9)$ is irrational; Fischler and Zudilin \cite[Theorems~4 and~6]{FZ10} refine Nesterenko's criterion by using
divisibility properties of the coefficients, and obtain for instance that $1,\zeta_q(i_0),\dots,\zeta_q(i_3)$ are
linearly independent for some odd $i_0<\dots<i_3$ with $i_0\le9$, $i_1\le37$, $i_2\le83$, $i_3\le145$. Neither
refinement changes the asymptotic constant, and we have not adapted either to units. In our setting the denominator of
$P_0$ divides $d_n^{A-1}$ in all cases we computed exactly (31 triples $(A,r,n)$ with $A\le12$), but we have not proved
this for $S_n$; if it held, $12A$ would become $12(A-1)$ in $\tau(A,r)$.
\end{remark}

% ===========================================================================
\section{Related work}\label{sec:related}
% ===========================================================================

% Source: handoff/fibzeta/many/NOVELTY_2.md (2026-11-19) and paper/NOVELTY_RECHECK.md. Re-run the Hexagon search
% (q = Lambert, Pisot, q-zeta; repos of twinheterodox) on the day of posting.
\emph{Even weights and weight one.} For even $k$ the values $\zeta_q(k)$ are polynomials in the Eisenstein series
$E_2,E_4,E_6$ at $q$, and at algebraic $q$ with $0<|q|<1$ the values of these series are algebraically independent by
Nesterenko's theorem \cite{Nes96}. At weight one there is a long line of irrationality and linear independence results
for Lambert series at algebraic bases: Tachiya \cite{Tac04} for algebraic integers of Pisot type, Bundschuh and
V\"a\"an\"anen \cite{BV05,BV05b}, Bundschuh \cite{Bun07} for integers of $\QQ$ or of an imaginary quadratic field, and
Dixit, Kumar and Pathak \cite{DKP24}, whose results on values of the $q$-exponential function at Pisot bases include the
weight-one value $\zeta_q(1)$. All of these concern weight at most two.

\emph{Odd weights at $q=1/p$.} For odd $k\ge3$ the known results concern $q=1/p$ with $p\in\ZZ\setminus\{0,\pm1\}$:
the dimension bound of Krattenthaler, Rivoal and Zudilin \cite{KRZ06}, with asymptotic constant
$\pi/(2\sqrt{\pi^2+12})=0.336\ldots$, and its refinements for fixed $A$ by Jouhet and Mosaki \cite{JM10} and Fischler
and Zudilin \cite{FZ10}; Postelmans and Van Assche \cite{PVA07} and Zudilin \cite{Zud02} treat weights one and two. The
extension of \cite{KRZ06} to algebraic $q$ sketched in \cite[\S3.5]{KRZ06} is conditional on a conjectured asymptotic for
the coefficients, and Bundschuh and Zudilin \cite[\S6]{BZ07} list the arithmetic of $\zeta_q(k)$ at reciprocals of
Pisot and Salem numbers as an open problem. The unrefereed draft \cite{UPL} claims $\zeta_{1/p}(3)\notin\QQ$ for integers
$p\ge2$, with a Lean formalisation of its main theorem (not checked by us), by a moving-denominator Hankel determinant.

\emph{Concurrent work.} The preprints \cite{Twi26,Twi26b}, discussed in the introduction, claim
$\zeta_{1/b}(3)\notin\QQ(b)$, and the linear independence of $1,\zeta_{1/b}(1),\zeta_{1/b}(2),\zeta_{1/b}(3)$ over
$\QQ(b)$, for the class of bases $b$ described in the introduction, which contains the Pisot and Salem numbers; the first
builds on \cite{UPL}. A third preprint by the same contributor \cite{Twi26c} extends the first claim to some
bases $b$ with several conjugates outside the unit disk, under a condition on the Mahler measure. These claims contain
Theorem~\ref{thm:main} as a special case, and the plain statement $\zeta_q(3)\notin\QQ$, which our method does not reach
(Remark~\ref{rem:plain}). Our proof of Theorem~\ref{thm:main} was found independently, uses a different construction
(the very-well-poised forms of \cite{KRZ06} and the reciprocity of Lemma~\ref{lem:recip}), and is checked in Lean. The
results of \cite{Twi26,Twi26b,Twi26c} stop at weight three; we know of no other lower bound for the dimension spanned by
odd $q$-zeta values at an algebraic irrational $q$ than the one of Section~\ref{sec:many}.

% ===========================================================================
\section{Formal verification}\label{sec:lean}
% ===========================================================================

The Lean development is the repository
\begin{center}
\url{https://github.com/telaviv/qzeta3-quadratic-units}
\end{center}
all file names below refer to it. Its Lean sources are those of commit \texttt{6318359}; later commits change only
documentation and this paper. It (library \lean{QZeta3}, Lean~4.35.0, Mathlib v4.35.0, about 3200 lines) proves
the following four theorems, stated in \lean{QZeta3/Final.lean}:
\begin{alltt}\footnotesize
theorem QZeta3.irrational_sqrt_mul_zetaq3 {\char123}q : \(\mathbb{R}\){\char125} {\char123}s : \(\mathbb{Z}\){\char125}
    (hq0 : q \(\ne\) 0) (hq1 : |q| < 1) (hs : q + q\(\sp{-1}\) = s)
    {\char123}d : \(\mathbb{N}\){\char125} (hd : 0 < d) (hsq : IsSquare ((d : \(\mathbb{Z}\)) * (s ^ 2 - 4))) :
    Irrational (Real.sqrt d * zetaq3 q)

theorem QZeta3.irrational_sub_inv_mul_zetaq3 {\char123}q : \(\mathbb{R}\){\char125} {\char123}s : \(\mathbb{Z}\){\char125}
    (hq0 : q \(\ne\) 0) (hq1 : |q| < 1) (hs : q + q\(\sp{-1}\) = s) :
    Irrational ((q - q\(\sp{-1}\)) * zetaq3 q)

theorem QZeta3.irrational_lucas_sum {\char123}a : \(\mathbb{N}\){\char125} {\char123}E : \(\mathbb{Z}\){\char125}
    (h : (E = 1 \(\land\) 1 \(\le\) a) \(\lor\) (E = -1 \(\land\) 3 \(\le\) a)) :
    Irrational (\(\sum\)' m : \(\mathbb{N}\), (m : \(\mathbb{R}\)) ^ 2 *
      ((lucasU a E (m + 1) : \(\mathbb{R}\)) / lucasU a E m - epsU a E))

theorem QZeta3.irrational_fib_sum :
    Irrational (\(\sum\)' m : \(\mathbb{N}\), (m : \(\mathbb{R}\)) ^ 2 *
      ((Nat.fib (m + 1) : \(\mathbb{R}\)) / Nat.fib m - Real.goldenRatio))
\end{alltt}
The definitions they use are, in \lean{QZeta3/Defs.lean},
\begin{alltt}\footnotesize
noncomputable def zetaq3 (q : \(\mathbb{R}\)) : \(\mathbb{R}\) :=
  \(\sum\)' m : \(\mathbb{N}\), (m : \(\mathbb{R}\)) ^ 2 * q ^ m / (1 - q ^ m)

def lucasU (a : \(\mathbb{N}\)) (E : \(\mathbb{Z}\)) : \(\mathbb{N}\) \(\to\) \(\mathbb{Z}\)
  | 0 => 0
  | 1 => 1
  | (m + 2) => a * lucasU a E (m + 1) + E * lucasU a E m

noncomputable def epsU (a : \(\mathbb{N}\)) (E : \(\mathbb{Z}\)) : \(\mathbb{R}\) :=
  (a + Real.sqrt ((a : \(\mathbb{R}\)) ^ 2 + 4 * E)) / 2
\end{alltt}
The proofs use only the axioms \lean{propext}, \lean{Classical.choice} and \lean{Quot.sound} (\lean{Axioms.lean});
there is no \lean{sorry}, no \lean{axiom} declaration, no \lean{native\_decide} and no \lean{implemented\_by}.
The definitions and the six ingredient statements, in the files \lean{QZeta3/Defs.lean} and
\lean{QZeta3/Statements.lean}, were fixed (``frozen'') in the first commit, before any ingredient was proved, and
were not changed afterwards. The four theorem statements above were committed with placeholder proofs before the last
two ingredients were proved, and their text is unchanged since then.

\emph{The statements.} In \lean{irrational\_sqrt\_mul\_zetaq3}, the hypotheses $q\ne0$, $|q|<1$ and $q+q^{-1}=s\in\ZZ$
say exactly that $q$ is a real quadratic unit of norm $+1$ with $0<|q|<1$ (as noted after Theorem~\ref{thm:main}), and
$d>0$ with $d(s^2-4)$ a square says exactly that $\QQ(\sqrt d)=\QQ(q)$; so it is Theorem~\ref{thm:main}.
\lean{irrational\_sub\_inv\_mul\_zetaq3} is the form with $\theta=q-q^{-1}$. In \lean{irrational\_lucas\_sum} the
recurrence is $U_{m+2}=aU_{m+1}+EU_m$, so $(a,E)=(P,-Q)$, and the hypothesis is admissibility; \lean{epsU a E} is
$\eta$. The $m=0$ terms of the series are $0$ in Lean (for \lean{zetaq3} the term is $0\cdot1/0=0$; in the Lucas
series the factor $m^2$ vanishes). \lean{irrational\_fib\_sum} is the case $a=E=1$, with Mathlib's \lean{Nat.fib} and
\lean{Real.goldenRatio} (\lean{QZeta3.lucasU\_one\_one}, \lean{QZeta3.epsU\_one\_one}).

\emph{What has to be trusted.} The kernel checks the proofs; to see that the statements mean what is claimed one reads
the four statements above, the three definitions, and Mathlib's \lean{tsum}, \lean{Irrational}, \lean{Real.sqrt},
\lean{IsSquare}, \lean{Nat.fib} and \lean{Real.goldenRatio}. A divergent \lean{tsum} has the value $0$, which is
rational, so none of the theorems can hold for a degenerate reason; in addition \lean{QZeta3.summable\_zetaq3} and
\lean{QZeta3.summable\_lucas} prove that the series converge. A blind back-translation of the statements, by a reviewer
(also a Claude model) who saw only the Lean source without comments or intended meaning, found no mismatch with the
author's working statement of Theorem~\ref{thm:main} and Corollaries~\ref{cor:lucas}--\ref{cor:fib} (the same content
as here); it is recorded in \lean{FIDELITY.md}. The
author also checked the statements against Theorem~\ref{thm:main} and the corollaries (\lean{STATEMENT\_REVIEW.md}),
and \lean{TRUST.md} summarises the trust base.

\subsection*{The chain}
\lean{QZeta3/Defs.lean} defines the construction for $A=4$, $r=1$ over an arbitrary field: the $q$-Pochhammer symbol
\lean{qPoch}, the rational function \lean{Rfun} of \eqref{eq:R}, the coefficients \lean{cc} of \eqref{eq:pf}, the
coefficients \lean{P3} and \lean{P0} of Lemma~\ref{lem:expansion} (with \lean{Wf} for $W_{s,j}$), and, over $\RR$,
\lean{dn} for $d_n$ and \lean{Sn} for $S_n$. \lean{QZeta3/Statements.lean} states six ingredients as propositions, and
\lean{QZeta3.irrational\_sub\_inv\_mul\_zetaq3\_of} (in \lean{QZeta3/Main.lean}) derives the theorem from them by the
argument of Section~\ref{sec:proof}; \lean{QZeta3/Reduction.lean} and \lean{QZeta3/Corollary.lean} derive the $\sqrt d$
form and the corollaries.
In the following table all Lean names are in the namespace \lean{QZeta3}.
\begin{center}
\footnotesize
\begin{tabular}{lll}
\toprule
Paper & Lean statement & Lean proof\\
\midrule
\eqref{eq:pf} & \lean{PFIdentity} & \lean{PartialFrac.pfIdentity}\\
Lemma~\ref{lem:expansion} & \lean{Expansion} & \lean{Expand.expansion}\\
Lemmas~\ref{lem:inv}, \ref{lem:recip}, Cor.~\ref{cor:cinv} & \lean{Recip} & \lean{Reciprocity.recip}\\
Lemma~\ref{lem:den}, \S\ref{sec:proof} & \lean{Integrality} & \lean{Denom.integrality}\\
Lemmas~\ref{lem:wp}, \ref{lem:size} & \lean{SizeBound} & \lean{Size.sizeBound}\\
Lemma~\ref{lem:cost} & \lean{CostBound} & \lean{Cost.costBound}\\
Theorem~\ref{thm:main}, $\theta=q-q^{-1}$ & & \lean{irrational\_sub\_inv\_mul\_zetaq3\_of}\\
Theorem~\ref{thm:main}, $\sqrt d$ & & \lean{irrational\_sqrt\_mul\_zetaq3\_ing}\\
Corollary~\ref{cor:lucas} & & \lean{lucas\_sum\_eq}\\
 & & \lean{irrational\_lucas\_sum\_ing}\\
\bottomrule
\end{tabular}
\end{center}
The formal statements are the lemmas in the form the assembly needs, for $A=4$, $r=1$ and real $q$. \lean{Expansion} and
\lean{Recip} are derived from \lean{PFIdentity}. \lean{Recip} states Lemma~\ref{lem:recip} at every real $q\ne0$ with
$|q|\ne1$. \lean{Integrality} says that $d_n^4P_0$ and $d_n^4P_3$ agree, at every real $x\ne0$ with $|x|\ne1$, with
one fixed integer Laurent polynomial (\lean{LaurentInt}); this is what allows evaluation at both $q$ and $q^{-1}$.
\lean{SizeBound} gives $S_n\ne0$ and $|S_n|\le C|q|^{n(n+1)}$ for large $n$, and \lean{CostBound} is
$|d_n(q)d_n(q^{-1})|^8\le C^n|q|^{-3(n+1)^2}$, the bound $\Tot(n)\le\frac38(n+1)^2$ of Lemma~\ref{lem:cost}.

\subsection*{Differences from the written proof}
\begin{itemize}
\item \emph{No number fields.} The formal proof does not use Mathlib's number fields. The role of $\sqrt d$ is played
by $\theta=q-q^{-1}$, as in Section~\ref{sec:proof}; $\sigma$ is replaced by evaluation of integer Laurent polynomials
at $q^{-1}$, and the norm by the product $p(q)\,p(q^{-1})$, which is an integer for $p\in\ZZ[X]$ because
$q+q^{-1}\in\ZZ$ (\lean{QZeta3.norm\_int}).
\item \emph{Partial-fraction coefficients.} The coefficients $c_{s,j}$ are \emph{defined} as the coefficients of the
power-series expansion of $u^4R_n$ at the pole $T=q^{-j}$, in the variable $u=1-q^jT$ (\lean{Gser}, \lean{cc}), and
\eqref{eq:pf} is then proved. With this definition, Corollary~\ref{cor:cinv} follows directly from Lemma~\ref{lem:inv},
without appeal to the uniqueness of partial fractions.
\item \emph{The cyclotomic bound.} The formal proof uses the elementary bound $\Tot(n)\le\frac38(n+1)^2$ of
Lemma~\ref{lem:cost}, with exponent $-\frac12$, not the asymptotic of Remark~\ref{rem:mertens}, which is not in Mathlib.
(The author's working notes used the asymptotic, with exponent $12/\pi^2-2$.)
\item \emph{Generality.} The formal statements fix $A=4$, $r=1$; Theorem~\ref{thm:main} needs nothing else. The Lean form
of Theorem~\ref{thm:main} allows every integer $d>0$ with $\QQ(\sqrt d)=\QQ(q)$, as does the statement above.
\end{itemize}

\subsection*{Provenance}
The library name \lean{QZeta3} (for the $q$-analogue of $\zeta(3)$) is also the name of the Lean formalisation
\cite{QZ3} of \cite{UPL}. The two are different developments, of different theorems, and ours is not a fork. One file,
\lean{QZeta3/Cost/Cyclo.lean} (cyclotomic size bounds), is adapted from the formalisation accompanying \cite{Kri26b},
which generalises that of \cite{Kri26a}, which in turn adapts part of \cite{QZ3} (Apache License~2.0); the file says so in
its header, and the repository's \lean{NOTICE} file records it. The mathematics of the present paper does not depend on
\cite{UPL}: the construction is that of \cite{KRZ06}.

\subsection*{Use of AI tools}
This work was carried out with substantial assistance from AI systems (Anthropic's Claude models), which were used to
search the literature, to develop and check the proof, to write the Lean formalisation and to draft this text. The
author directed the project and is responsible for its content. The formal proof is checked by the Lean kernel and
uses only the standard axioms \lean{propext}, \lean{Classical.choice} and \lean{Quot.sound}.

% ===========================================================================
\bibliographystyle{amsplain}
\bibliography{refs}

\end{document}
